
\section{Literature Review}
\label{sec:litreview}

Decentralized databases built on peer-to-peer (P2P) storage networks, such as the InterPlanetary File System (IPFS), offer resistance to censorship, data sovereignty, and fault tolerance, but come at a steep price in query performance. Every data access requires one or more P2P network fetches -- each orders of magnitude slower than a local disk read, making network I/O the dominant bottleneck~\cite{trautwein2022design, henningsen2020mapping, daniel2022ipfs}. Although decades of research have optimized cache and indexing for centralized and cloud-hosted databases~\cite{graefe2011modern, armbrust2021lakehouse}, these techniques assume low-latency local storage and do not apply to decentralized settings where queries must traverse Merkle DAGs over content-addressed P2P networks. 
% This review surveys work that informs \ours{}, which contributes (1)~a TEE-based, multi-layer semantic-aware caching system with SMT-based query containment, and (2)~a content-addressed, multi-dimensional sharding technique that enables predicate pushdown at the IPFS storage layer.


%% ---------------------------------------------------------------
\subsection{IPFS and Decentralized Databases}
\label{sec:ipfs}
%% ---------------------------------------------------------------

IPFS~\cite{benet2014ipfs} introduced content-addressed storage on a Kademlia-based DHT with Merkle DAG data structures, where every object is identified by its cryptographic Content Identifier (CID). Large-scale measurements report that the median IPFS retrieval latency is 2-5 seconds for unpinned content, with multi-hop DHT lookups adding hundreds of milliseconds per round~\cite{trautwein2022design, henningsen2020mapping, daniel2022ipfs}. While IPLD~\cite{ipld2023spec} provides a linked-data model layer atop IPFS, it doesn't natively support pushdown storage-level filtering for relational data to reduce I/O costs.

Several systems build database abstractions on decentralized storage. Early efforts, such as BigChainDB~\cite{mcconaghy2016bigchaindb}, overlay database semantics on blockchain consensus but inherit throughput limitations, while key-value and document interfaces on IPFS lack relational query support. Minerva~\cite{yao2024minerva} distributes query execution across IPFS peers with a cost-based planner but does not incorporate caching or predicate-aware data organization, leaving content retrieval latency unaddressed.

Most relevant to our work are MtDB~\cite{hossain2025mtdb} and Web3DB~\cite{mukherjee2025web3db}. MtDB integrates IPFS storage, Ethereum smart contracts for metadata coordination, and Intel SGX for secure query execution; its delta-based on-chain index enables partition discovery but not predicate pushdown within a partition. Web3DB provides blockchain-managed access control for relational queries, but lacks enclave-based secure query execution. Both of these systems lack a caching or data organization layer to reduce I/O bottlenecks. \ours{} builds on MtDB's architecture, adding semantic caching and content-addressed sharding to reduce IPFS I/O and query latency.


%% ---------------------------------------------------------------
\subsection{Database Caching and Semantic Caching}
\label{sec:caching}
%% ---------------------------------------------------------------

Traditional database caching -- buffer pools~\cite{graefe2011modern}, key-value stores such as Memcached~\cite{fitzpatrick2004distributed}, Redis~\cite{redis2009} and result caches in cloud warehouses~\cite{dageville2016snowflake, melnik2010dremel} -- reduces I/O to local or cloud-attached disk, where cache miss penalties are microseconds. In distributed settings, DistCache~\cite{liu2019distcache} provides provably balanced caching for large-scale storage, Crystal~\cite{durner2021crystal} unifies columnar and row-based caching for analytics, and the work of Tang et al.~\cite{tang2024data} coordinates multi-layer caching for petabyte-scale OLAP. Metadata caching in Presto~\cite{wang2022metadata} eliminates redundant metastore calls, while HyCache~\cite{zhao2013hycache} interposes application-aware caching in distributed file systems. All these systems target predictable, low-latency storage and do not address the fundamentally different cost model of P2P content retrieval, where a cache miss costs 100-5000 ms.

Result caches keyed by normalized SQL strings~\cite{dageville2016snowflake, melnik2010dremel} cannot exploit overlap between queries: a query for \texttt{Age > 60} is a miss even if \texttt{Age > 50} is cached. Semantic caching addresses this by annotating cached results with defining predicates and checking containment~\cite{dar1996semantic}. Keller and Basu~\cite{keller1996predicate} extended the approach with predicate indices, and Roy et al.~\cite{roy2000don} advocated caching intermediate results for cross-query reuse. Most recently, You et al. \cite{you2025soc} propose succinct semantic representations for OLAP caching but rely on heuristic predicate overlap rather than formal containment, and target cloud warehouses with fast storage. We use SMT-based containment via Z3~\cite{de2008z3} to guarantee soundness, and the extreme IPFS miss penalty makes the 2-15 ms predicate-checking overhead negligible.

Cache eviction is a core component of any caching system -- especially in an in-memory cache, given limited storage capacity. For eviction, classical policies \cite{megiddo2003arc, einziger2017tinylfu,jiang2002lirs,zhang2024sieve} treat each cache object as opaque. Recently, the SIEVE algorithm~\cite{zhang2024sieve} achieves both eager promotion and quick demotion of cache entries using a single visited bit, but it ignores semantic relationships, which are essential in our scenario. Cost-aware eviction consistently outperforms cost-oblivious policies~\cite{podlipnig2003survey} when cache miss cost is high (e.g., 100--5000 ms). Our cache eviction algorithm augments SIEVE with a second utility bit encoding the semantic reuse potential of each entry, preserving O(1) complexity.


%% ---------------------------------------------------------------
\subsection{Query Containment and Predicate Pushdown}
\label{sec:containment}
%% ---------------------------------------------------------------

Query containment -- determining whether one query's results are a subset of another's -- was formalized by Chandra and Merlin~\cite{chandra1977optimal} as NP-complete for conjunctive queries. SMT solvers~\cite{de2008z3} provide a natural encoding for containment checks. VeriEQL~\cite{he2024verieql} leverages Z3 for bounded equivalence verification of SQL queries under integrity constraints, while Cosette~\cite{chu2017cosette} approaches query equivalence via symbolic constraint solving. These systems are designed for offline verification; in contrast, our containment checker operates online, where soundness (no false positives) is required but incompleteness is acceptable.

Predicate pushdown is a standard optimization in lakehouse and columnar systems~\cite{armbrust2021lakehouse}. FlexPushdownDB~\cite{yang2021flexpushdowndb} dynamically decides whether to push predicates to cloud storage (e.g., S3 Select) or process them locally, trading off network cost against computation. In decentralized databases~\cite{hossain2025mtdb, mukherjee2025web3db}, however, the pushdown target is not a cloud service with native filtering support, but an IPFS node that exposes no query interface. We address this gap by constructing a custom DAG over IPFS that encodes range metadata using IPLD~\cite{ipld2023spec}, enabling predicate filtering at the storage layer without any intermediate middleware. To our knowledge, no prior work has achieved pushdown filtering directly at the IPFS level.


%% ---------------------------------------------------------------
\subsection{Data Sharding and Indexing}
\label{sec:sharding}
%% ---------------------------------------------------------------

Range, hash, and composite partitioning are well-established for parallel databases~\cite{dewitt1992parallel}. Cloud-native systems like CockroachDB~\cite{taft2020cockroachdb} and Spanner~\cite{bacon2017spanner} use adaptive range partitioning with automatic shard splitting. DynaHash~\cite{luo2022dynahash} addresses efficient data rebalancing in AsterixDB via extensible hashing. All assume mutable, low-latency storage; IPFS's immutable content-addressed model requires constructing new Merkle DAGs rather than rebalancing in place.

Indexing for decentralized storage remains limited. Existing IPFS indexing frameworks target key-value or keyword lookups~\cite{arer2022efficient}, and blockchain-anchored indices focus on provenance~\cite{rathee2019blockchain} rather than query optimization. The Merkle DAG~\cite{benet2014ipfs, merkle1987digital} is a natural index for content-addressed systems. Our technique constructs a custom DAG in which internal nodes record min/max ranges across multiple attributes, enabling top-down traversal with branch pruning -- analogous to B$^+$-tree scans but adapted to immutable IPFS. Learned indices~\cite{kraska2018case} and related adaptive indexing techniques assume mutable storage and do not transfer to our setting.


%% ---------------------------------------------------------------
\subsection{TEE and Blockchain-based Databases}
\label{sec:tee}
%% ---------------------------------------------------------------

Trusted execution environments (TEEs) provide isolated execution with protected memory and confidentiality guarantees. Intel SGX~\cite{costan2016intel}, one of the most widely used TEEs, supports application-level enclaves and remote attestation. TEEs are well-suited for secure database caching, where both the cache state and the query execution logic can be carefully scoped to fit within enclave memory.

Several database systems have leveraged TEEs for secure query processing. EnclaveDB~\cite{priebe2018enclavedb} executes SQL operators inside SGX enclaves, while Opaque~\cite{zheng2017opaque} integrates SGX with Spark to support secure data analytics. CryptDB~\cite{popa2011cryptdb}, although not TEE-based, enables SQL queries over encrypted data using adjustable encryption schemes, but still leaks access patterns. However, these systems are designed for local or cloud-backed storage environments with relatively low-latency data access. As a result, they do not address the core challenges of peer-to-peer storage latency.

A related line of work combines blockchain systems with confidential computation. For example, Enigma~\cite{zyskind2018enigma} uses blockchain coordination together with privacy-preserving computation, but its primary goal is computation confidentiality rather than query acceleration. Similarly, verifiable query processing systems such as vChain~\cite{xu2019vchain} focus on correctness guarantees through cryptographic proofs, often incurring additional overhead without improving latency. MtDB~\cite{hossain2025mtdb} provides the underlying blockchain-based metadata coordination model that we build upon. In contrast, \ours{} introduces new mechanisms entirely at the off-chain caching and data-organization layers, targeting lower query latency and higher efficiency for decentralized database workloads.


% %% ---------------------------------------------------------------
% \subsection{Research Gaps}
% \label{sec:gaps}
% %% ---------------------------------------------------------------

% Three gaps emerge from the literature. \textbf{(1)~No semantic caching for decentralized databases}: semantic caching~\cite{dar1996semantic, keller1996predicate, you2025soc} has been studied only in centralized settings, despite the extreme miss penalty of IPFS (100--5000~ms vs.\ $<$1~ms for local disk) making it essential rather than optional. \textbf{(2)~No predicate pushdown at the content-addressed storage layer}: existing IPFS indexing~\cite{li2024ipfs, rathee2019blockchain, arer2022efficient} supports only key-value or keyword lookups and cannot push SQL predicates into content retrieval. \textbf{(3)~No TEE-optimized semantic eviction}: current policies~\cite{zhang2024sieve, megiddo2003arc} treat cached objects as independent, ignoring containment relationships where a broad query's result can serve many narrower queries. These gaps are interconnected: semantic caching reduces the number of IPFS fetches, content-addressed sharding reduces the cost per fetch, and TEE-optimized eviction maximizes cache utility within limited enclave memory.


%% ========================================================================
%% REFERENCES
%% ========================================================================

% \clearpage
% \setcounter{page}{1}


\bibliographystyle{abbrv}
\bibliography{references}


% \setcounter{page}{1}


% \begin{center}
% \textbf{\large Appendix}
% \end{center}

% \subsection*{Other data or stuff}

% The proposal is limited to 15 pages, but you can add extra stuff here if you need.